% File: Volume-III-Extensions/appendices/appJ-category-theory.tex

\chapter{Category Theory Reference}
\label{appJ-category-theory}

This appendix collects a compact set of categorical notions used throughout
Volume~III, with an emphasis on those aspects most relevant to homotopy
theoretic formulations of geometric computation, SpherePop calculus,
and lamphrodynamic semantic evolution.  It is not intended as a complete
introduction to category theory, but as a technical reference summarizing
definitions, constructions, and theorems that are explicitly invoked in the
main text.

We assume basic familiarity with ordinary categories and functors.  Throughout,
$\Cat$ denotes the $2$--category of (small) categories, and $\Set$ denotes the
category of sets.  When needed, we indicate size issues by assuming the
existence of universes, although these subtleties do not affect the results
presented here.

%------------------------------------------------------------------
\section{Categories and Functors}
%------------------------------------------------------------------

\begin{definition}
A \emph{category} $\C$ consists of a class of objects $\Ob(\C)$, a class of
morphisms $\Hom_\C(X,Y)$ for each pair of objects $(X,Y)$, identity morphisms
$\id_X \in \Hom_\C(X,X)$, and a composition law
$\Hom_\C(X,Y)\times\Hom_\C(Y,Z)\to\Hom_\C(X,Z)$
that is associative and unital.
\end{definition}

\begin{definition}
A \emph{functor} $F:\C\to\D$ assigns to each object $X\in\C$ an object $F(X)\in\D$
and to each morphism $f:X\to Y$ a morphism $F(f):F(X)\to F(Y)$, preserving
identities and composition.
\end{definition}

\begin{definition}
A \emph{natural transformation} $\eta:F\Rightarrow G$ between functors
$F,G:\C\to\D$ assigns to each $X\in\C$ a morphism
$\eta_X: F(X)\to G(X)$ such that for every $f:X\to Y$ in $\C$, the naturality
square $G(f)\circ\eta_X = \eta_Y\circ F(f)$ commutes.
\end{definition}

%------------------------------------------------------------------
\section{Limits and Colimits}
%------------------------------------------------------------------

Limits and colimits abstract universal constructions.  These are essential in
SpherePop, where merge operations correspond to homotopy colimits of geometric
data, and collapse operations correspond to truncation functors.

\begin{definition}
Given a diagram $D:J\to\C$, a \emph{limit} $\varprojlim D$ is a universal cone
over $D$.  A \emph{colimit} $\varinjlim D$ is a universal cocone.
\end{definition}

\begin{example}
Products, equalizers, pullbacks are special cases of limits; coproducts,
coequalizers, and pushouts are colimits.
\end{example}

\begin{remark}
In homotopy theoretic settings (model categories or $\infty$--categories)
one replaces limits and colimits with \emph{homotopy} limits and colimits,
which rectify the failure of strict universal properties under weak
equivalences.  This distinction is central to the formal proof that SpherePop
merge operations realize homotopy colimits rather than strict colimits.
\end{remark}

%------------------------------------------------------------------
\section{Adjunctions}
%------------------------------------------------------------------

\begin{definition}
An \emph{adjunction} between categories $\C$ and $\D$ is a pair of functors
$F:\C\to\D$ and $G:\D\to\C$ together with natural bijections
\[
\Hom_\D(FX,Y)\cong \Hom_\C(X,GY),
\]
natural in $X\in\C$ and $Y\in\D$.  We write $F\dashv G$.
\end{definition}

Adjunctions generate many fundamental constructions in homotopy theory,
including derived functors and Quillen adjunctions.

%------------------------------------------------------------------
\section{Model Categories}
%------------------------------------------------------------------

To work with explicit homotopy theoretic constructions (beyond derived
$\infty$--categorical formalism), the notion of a model category is often
useful.

\begin{definition}
A \emph{model category} $(\M,\W,\Fib,\Cof)$ consists of a category $\M$ together
with three classes of morphisms: weak equivalences $\W$, fibrations $\Fib$, and
cofibrations $\Cof$, satisfying the Quillen axioms.
\end{definition}

Every model category admits derived functors and homotopy limits/colimits via
cofibrant and fibrant replacements.  In Volume~III we use model structures
primarily as a conceptual bridge to $\infty$--categorical homotopy theory,
rather than as a computational tool.

%------------------------------------------------------------------
\section{$\infty$--Categories}
%------------------------------------------------------------------

\begin{definition}
An $\infty$--category is a higher category whose $n$--morphisms are all
invertible for $n>1$.  Many equivalent models exist: quasi--categories,
complete Segal spaces, simplicial categories, etc.
\end{definition}

\begin{remark}
For most of this volume, the intended model is that of quasi--categories, as in
Lurie.  Homotopy limits and colimits are encoded by mapping spaces rather than
strict diagrams.
\end{remark}

\begin{definition}
A functor between $\infty$--categories preserves equivalences and coherent
higher morphisms.  Natural transformations generalize to higher homotopies
between such functors.
\end{definition}

%------------------------------------------------------------------
\section{Derived Functors}
%------------------------------------------------------------------

Let $F:\C\to\D$ be a functor between model categories.  Its \emph{left derived
functor} $\mathbf{L}F$ is obtained by applying $F$ to a cofibrant replacement;
the \emph{right derived functor} $\mathbf{R}F$ uses fibrant replacement.

\begin{remark}
In $\infty$--categorical settings, derived functors are simply the homotopy
correct functors defined on $\infty$--categories themselves; no replacements
are needed if one works natively in quasi--categories.
\end{remark}

%------------------------------------------------------------------
\section{Homotopy Colimits}
%------------------------------------------------------------------

\begin{definition}
Given a diagram $D:J\to\M$ in a model category, the \emph{homotopy colimit}
$\hocolim D$ is the left derived colimit, typically computed via cofibrant
replacement or via simplicial resolutions.
\end{definition}

\begin{remark}
In Volume~III, Theorem~24.2 identifies SpherePop merge operations with homotopy
colimits in a suitable $\infty$--categorical framework.  This justifies the
intuition that merge is a higher--categorical ``gluing'' operation preserving
homotopical information.
\end{remark}

%------------------------------------------------------------------
\section{Quillen Adjunctions and Equivalences}
%------------------------------------------------------------------

\begin{definition}
A \emph{Quillen adjunction} between model categories $\M$ and $\N$ is an
adjunction $F\dashv G$ such that $F$ preserves cofibrations and acyclic
cofibrations (equivalently, $G$ preserves fibrations and acyclic fibrations).
\end{definition}

\begin{definition}
A Quillen adjunction $F\dashv G$ is a \emph{Quillen equivalence} if the derived
unit and counit maps are weak equivalences.  In this case, $\M$ and $\N$ have
equivalent homotopy categories.
\end{definition}

\begin{remark}
Quillen equivalence is the key mechanism by which SpherePop configurations are
shown to be \emph{computationally invariant} under change of model, thereby
allowing derived truncation and homotopy colimit constructions to be interpreted
in multiple equivalent frameworks.
\end{remark}

%------------------------------------------------------------------
\section{Truncation and Postnikov Towers}
%------------------------------------------------------------------

\begin{definition}
For an $\infty$--category $\C$, the $n$--truncation $\tau_{\le n}\C$ is obtained
by forcing all homotopy groups above degree $n$ to vanish.  At the level of
objects, this corresponds to retaining data up to homotopies of order $n$.
\end{definition}

\begin{remark}
Collapse operations in SpherePop are conjecturally equivalent to derived
truncation functors, eliminating higher--order homotopic detail while
preserving essential information.  Conjecture~20.3 provides a precise
formulation and evidence for this interpretation.
\end{remark}

%------------------------------------------------------------------
\section{Monoidal Structures}
%------------------------------------------------------------------

Many constructions in RSVP and SpherePop rely on monoidal structures, such as
tensor products of complexes and monoidal $\infty$--categories.

\begin{definition}
A \emph{monoidal category} $(\C,\otimes,\mathbf{1})$ is a category equipped with
a bifunctor $\otimes:\C\times\C\to\C$, associativity constraints, and a unit
object $\mathbf{1}$ satisfying coherence conditions.
\end{definition}

\begin{remark}
The monoidal structure on derived categories of sheaves underlies the AKSZ
construction; similarly, monoidal structures on $\infty$--categories support
the compositional semantics of SpherePop computations.
\end{remark}

%------------------------------------------------------------------
\section*{References}
%------------------------------------------------------------------

For further reading:
\begin{itemize}
\item Lurie: \emph{Higher Topos Theory}, \emph{Higher Algebra}.
\item Hovey: \emph{Model Categories}.
\item Dwyer--Hirschhorn--Kan--Smith: \emph{Homotopy Theories and Model Categories}.
\item Riehl: \emph{Category Theory in Context}.
\end{itemize}


